% SM Table S20: Top-10 Polypharmacological Candidates with Retrosynthetic Analysis
% Generated by r8b_pipeline.py (ASKCOS public API)

\begin{table}[htbp]
\centering
\caption{Top-10 polypharmacological candidates from the generative library. MPO $\geq$ 0.50 across multiple targets. \label{tab:sm_s20_top10_retrosynthesis}}
\small
\begin{tabular}{l l l l l l}
\toprule
Rank & SMILES & MPO Score & SYBA Score & QED & Targets \\
\midrule
1 & Cc1ccc(CN2CCN(Cc3ccccc3)CC2)cc1O & 0.829 & 109.0 & 0.939 & PfDHFR, PfCRT, PfATP4 \\
2 & Cc1ccccc1CN1CCN(Cc2cccc(O)c2)CC1 & 0.828 & 124.9 & 0.939 & PfDHFR, PfCRT, PfATP4 \\
3 & Cc1cc(CN2CCN(Cc3ccccc3)CC2)ccc1O & 0.826 & 100.6 & 0.939 & PfDHFR, PfCRT, PfATP4 \\
4 & COc1ccc(CN2CCN(Cc3ccccc3C)CC2)cc1O & 0.826 & 133.7 & 0.916 & PfDHFR, PfCRT, PfATP4 \\
5 & Oc1cccc(CN2CCN(Cc3ccc(Cl)cc3)CC2)c1 & 0.826 & 116.0 & 0.937 & PfDHFR, PfCRT, PfATP4 \\
6 & Cc1ccc(CN2CCN(Cc3ccc(O)cc3)CC2)c(C)c1 & 0.826 & 122.0 & 0.938 & PfDHFR, PfCRT, PfATP4 \\
7 & COc1ccc(CN2CCN(Cc3ccc(C)c(O)c3)CC2)cc1 & 0.826 & 119.1 & 0.916 & PfDHFR, PfCRT, PfATP4 \\
8 & Cc1ccc(CN2CCN(Cc3cccc(O)c3)CC2)c(C)c1 & 0.825 & 129.5 & 0.938 & PfDHFR, PfCRT, PfATP4 \\
9 & Oc1ccc(CN2CCN(Cc3ccccc3)CC2)cc1Cl & 0.825 & 113.8 & 0.937 & PfDHFR, PfCRT, PfATP4 \\
10 & Oc1ccc(CN2CCN(Cc3ccccc3Cl)CC2)cc1 & 0.824 & 122.5 & 0.937 & PfDHFR, PfCRT, PfATP4 \\
\bottomrule
\end{tabular}
\end{table}

\begin{table}[htbp]
\centering
\caption{Predicted retrosynthetic routes for top-10 candidates via ASKCOS (public API, MIT). Precursors and reaction types from template relevance model (Reaxys database).
\label{tab:sm_s20_retrosynthesis}}
\small
\begin{tabular}{l l l}
\toprule
Compound & Proposed Precursors & Reaction Type \\
\midrule
T1 (route 1) & Cc1ccc(C=O)cc1O.c1ccc(CN2CCNCC2)cc1 & Reductive amination (aldehyde + amine) \\
T1 (route 2) & C1CNCCN1.C=O.C=O.Cc1ccccc1O.c1ccccc1 & Formylation / multicomponent \\
T1 (route 3) & COc1cc(CN2CCN(Cc3ccccc3)CC2)ccc1C & O-alkylation / ether formation \\
T2 (route 1) & Cc1ccccc1CN1CCNCC1.O=Cc1cccc(O)c1 & Reductive amination (aldehyde + amine) \\
T2 (route 2) & Cc1ccccc1C=O.Oc1cccc(CN2CCNCC2)c1 & Reductive amination (aldehyde + amine) \\
T2 (route 3) & Cc1ccccc1CN1CCNCC1.Oc1cccc(CBr)c1 & N-alkylation (alkyl bromide) \\
T3 (route 1) & Cc1cc(C=O)ccc1O.c1ccc(CN2CCNCC2)cc1 & Reductive amination (aldehyde + amine) \\
T3 (route 2) & Cc1ccccc1O.c1ccc(CN2CCNCC2)cc1 & C-N cross-coupling \\
T3 (route 3) & C1CNCCN1.Cc1cc(C=O)ccc1O.O=Cc1ccccc1 & Reductive amination (two aldehydes + amine) \\
T4 (route 1) & COc1ccc(C=O)cc1O.Cc1ccccc1CN1CCNCC1 & Reductive amination (aldehyde + amine) \\
T4 (route 2) & C1CNCCN1.C=O.C=O.COc1ccccc1O.Cc1ccccc1 & Formylation / multicomponent \\
T4 (route 3) & COc1ccc(CO)cc1O.Cc1ccccc1CN1CCNCC1 & N-alkylation (alcohol) \\
T5 (route 1) & Clc1ccc(CN2CCNCC2)cc1.O=Cc1cccc(O)c1 & Reductive amination (aldehyde + amine) \\
T5 (route 2) & O=Cc1ccc(Cl)cc1.Oc1cccc(CN2CCNCC2)c1 & Reductive amination (aldehyde + amine) \\
T5 (route 3) & Clc1ccc(CN2CCNCC2)cc1.Oc1cccc(CBr)c1 & N-alkylation (alkyl bromide) \\
T6 (route 1) & Cc1ccc(C=O)c(C)c1.Oc1ccc(CN2CCNCC2)cc1 & Reductive amination (aldehyde + amine) \\
T6 (route 2) & C1CNCCN1.Cc1ccc(C=O)c(C)c1.O=Cc1ccc(O)cc1 & Reductive amination (two aldehydes + amine) \\
T6 (route 3) & C1CNCCN1.Cc1ccc(CCl)c(C)c1.Oc1ccc(CCl)cc1 & N-alkylation (alkyl chloride) \\
T7 (route 1) & COc1ccc(CN2CCNCC2)cc1.Cc1ccc(C=O)cc1O & Reductive amination (aldehyde + amine) \\
T7 (route 2) & COc1ccc(C=O)cc1.Cc1ccc(CN2CCNCC2)cc1O & Reductive amination (aldehyde + amine) \\
T7 (route 3) & COc1ccc(CN2CCNCC2)cc1.Cc1ccc(CCl)cc1O & N-alkylation (alkyl chloride) \\
T8 (route 1) & Cc1ccc(C=O)c(C)c1.Oc1cccc(CN2CCNCC2)c1 & Reductive amination (aldehyde + amine) \\
T8 (route 2) & Cc1ccc(CCl)c(C)c1.Oc1cccc(CN2CCNCC2)c1 & N-alkylation (alkyl chloride) \\
T8 (route 3) & Cc1ccc(CBr)c(C)c1.Oc1cccc(CN2CCNCC2)c1 & N-alkylation (alkyl bromide) \\
T9 (route 1) & O=Cc1ccc(O)c(Cl)c1.c1ccc(CN2CCNCC2)cc1 & Reductive amination (aldehyde + amine) \\
T9 (route 2) & Oc1ccccc1Cl.c1ccc(CN2CCNCC2)cc1 & C-N cross-coupling \\
T9 (route 3) & C1CNCCN1.C=O.C=O.Oc1ccccc1Cl.c1ccccc1 & Formylation / multicomponent \\
T10 (route 1) & Clc1ccccc1CN1CCNCC1.O=Cc1ccc(O)cc1 & Reductive amination (aldehyde + amine) \\
T10 (route 2) & O=Cc1ccccc1Cl.Oc1ccc(CN2CCNCC2)cc1 & Reductive amination (aldehyde + amine) \\
T10 (route 3) & C1CNCCN1.O=Cc1ccc(O)cc1.O=Cc1ccccc1Cl & Reductive amination (two aldehydes + amine) \\
\bottomrule
\end{tabular}
\end{table}
